\subsection{Limits of Functions}


Let $D \subseteq \mathbb{R}$ and $f : D \rightarrow \mathbb{R}$.

\begin{definition}
  Let $S \subseteq \mathbb{R}$ and $c \in \mathbb{R}$. We say $c$ is a {\bf cluster point} of $S$ if there exists $(x_n)_{n \in \mathbb{N}}$ in $S \setminus \{c\}$ such that $\lim_{n \rightarrow \infty} x_n = c$.
\end{definition}
\begin{definition}
  Let $D \subseteq \mathbb{R}$ and $f : D \rightarrow \mathbb{R}$. Let $L \in \mathbb{R}$ and let $c$ be a cluster point of $D$.

  We say { $f$ \bf converges to $L$} as $x$ tends to $c$, or $\lim_{x \rightarrow c} f(x) = L$ if
  \[\forall \epsilon >0 \quad \exists \delta > 0 \quad \forall x \in D \quad \text{s.t.}\]
  \[ \text{ if } 0 < | x - c| < \delta \text{ then } |f(x) - L | < \epsilon\]
\end{definition}
$ $ \\
In the above definition, the definition of $f(c)$ is irrelevent.

\begin{theorem}
  Let $c \in \mathbb{R}$ be a cluster point of $D \subseteq \mathbb{R}$. Let $f: D \rightarrow \mathbb{R}$ and $L \in \mathbb{R}$. Then the following are equivalent:
  \begin{itemize}
    \item $\lim_{x \rightarrow c} f(x) = L$
    \item For any sequence $(x_n)_{n \in \mathbb{N}}$ in $D \setminus \{c\}$ 
      \[\lim_{n \rightarrow \infty} x_n = c \implies \lim_{n\rightarrow \infty} f(x_n) = L.\]
  \end{itemize}
\end{theorem}
\begin{proof}
    \begin{itemize}
      \item Suppose $\lim_{x \rightarrow c} f(x) = L$ and $\lim_{n \rightarrow \infty} x_n = c$. Now fix $\epsilon > 0$. Then we obtain a $\delta$ where
        \[0 < |x - c| < \delta \implies |f(x) - L| < \epsilon.\]
        Using this $\delta$, we obtain an $N$ such that
        \[\forall n > N \quad |x_n - c| < \delta.\]
        Then since $x_n \not = c$ for all $x$, the conditions for $f(x_n) \rightarrow L$ are met.
      \item Since assuming an implication is difficult to work with, we instead prove the contrapositive: 
        \begin{align*}
        \lim_{x \rightarrow c} f(x) \not = L \wedge x_n \rightarrow c \implies \forall \epsilon > 0, \  \exists N, \ \forall n \geq N \\
        |f(x_n) - L| > \epsilon
        .\end{align*}
        The negation of the limit of $f$ gives us an $\epsilon$ where for any $\delta$ we can find an $x$ such that $0 < |x - c| < \delta$ and $| f(x) - L| > \epsilon$. Then, using this $\epsilon$, fix $\delta >0$. Choose $N$ such that
        \[| x - c| < \delta\]
        and then the conclusion follows.
    \end{itemize}
\end{proof}

\begin{prop}
  Let $c$ be a cluster point of $D \subseteq \mathbb{R}$. Suppose there exists a sequence $(x_n)_{n\in \mathbb{N}}$ and $(y_n)_{n \in \mathbb{N}}$ in $D \setminus \{c\}$ such that
  \begin{itemize}
    \item $\lim_{n \rightarrow \infty} x_n = c$ and $\lim_{y \rightarrow \infty} = c$
    \item $\lim_{n\rightarrow \infty} f(x_n) = A$ and $\lim_{n \rightarrow \infty} f(y_n)= B$ with $A \not = B$.
  \end{itemize}

  Then we conclude $\lim_{x \rightarrow c} f(x)$ cannot exist.
\end{prop}
\begin{proof}
  Suppose the limit exists. Then since $x_n$ and $y_n$ have the same limit
  \[\lim_{n \rightarrow \infty} f(x_n) = A = B = \lim_{n \rightarrow \infty} f(y_n)\]
  which is a contradiction.
\end{proof}

\begin{prop}
    Let $c$ be a cluster point of $D \subseteq \mathbb{R}$ and let $f : D \rightarrow \mathbb{R}$.

    If $\lim_{x \rightarrow c} f(x) = L_1$ and $\lim_{x \rightarrow c} f(x) = L_2$ then
    \[L_1 = L_2.\]
\end{prop}
\begin{proof}
   $\lim_{x \rightarrow c} f(x) = L$ is equivalent to $f(x_n) \rightarrow L$ for $x_n \rightarrow c$. Then
   \[\lim_{n \rightarrow \infty} f(x_n) = L_1 = L_2.\]
   So $L_1 = L_2$ by uniqueness of limits.
\end{proof}

\begin{theorem} (Algebra of limits for functions)

  Suppose $D \subseteq \mathbb{R}$ and $c$ is a cluster point of $D.$ Let $f , g: D \rightarrow \mathbb{R}$ be two functions.

  Suppose $\lim_{x \rightarrow c} f(x) = L$ and $\lim_{x \rightarrow c} g(x) = M.$

  Let $x \in \mathbb{R}$. Then
  \begin{itemize}
    \item $\lim_{x \rightarrow c}  \ \left(f(x) + g(x)\right) = L + M$
    \item $\lim_{x \rightarrow c} \ \alpha f(x) = \alpha L$
    \item $\lim_{x \rightarrow c} \ f(x) g(x) = LM$
    \item If $M \not = 0$, then $\lim_{x \rightarrow c} \ \frac{f(x)}{g(x)} = \frac{L}{M}$
  \end{itemize}
\end{theorem}
\begin{proof}
    By turning the limit of the function into a sequence limit, they can be easily proved using the algebra of limits.
\end{proof}

\begin{theorem}
  Suppose $D \subseteq \mathbb{R}$, $c \in \mathbb{R}$ is a cluster point of $D$, $f : D \rightarrow \mathbb{R}$ and $\lim_{x \rightarrow c} f(x) = L$. Then
  \begin{itemize}
    \item If $f(x) \leq M \quad \forall x \in D$, then $L \leq M$
    \item If $f(x) \geq M \quad \forall x \in D$, then $L \geq M$.
  \end{itemize}
\end{theorem}
\begin{proof}
    Consider WLOG $f(x) \leq M$. Then for $x_n \rightarrow \infty$
    \[\lim_{n \rightarrow \infty} f(x_n) \leq M\]
    so since limits preserve non-strict inequalities, $L \leq M$.
\end{proof}

\begin{theorem}
  Suppose $D \subseteq \mathbb{R}$, $c \in \mathbb{R}$ is a cluster point of $D$ and $f,g,h : D \rightarrow \mathbb{R}$ with
  \[f(x) \leq g(x) \leq h(x) \quad \forall x \in D.\]
  and
  \[\lim_{x \rightarrow c} f(x) = L = \lim_{x \rightarrow c} h(x).\]
  Then $\lim_{x \rightarrow c} g(x) = L$.
\end{theorem}
\begin{proof}
    Consider $x_n \rightarrow c$. Then
    \begin{align*}
      f(x_n) &\leq g(x_n) \leq h(x_n) \\
      \lim_{n \rightarrow \infty} f(x_n) &\leq \lim_{n \rightarrow \infty} g(x_n) \leq \lim_{n \rightarrow \infty} h(x_n)
    .\end{align*}
    so by sandwhich theorem the limit of $g$ at $c$ is $L$.
\end{proof}

Now we will consider the limit of a function to $\infty$.
\begin{definition}
    Let $D \subseteq \mathbb{R}$, $f : D \rightarrow \mathbb{R}$ and $L \in \mathbb{R}$. Suppose there exists a sequence $(a_n)_{n \in \mathbb{N}}$ in $D$ such that $\lim_{n \rightarrow \infty} a_n = \infty$.

    We say $f$ converges to $L$ as $x \rightarrow \infty$ if
    \[\forall \epsilon > 0 \quad \exists N \in \mathbb{R} \text{    s.t.  } \forall x \in D, \quad x > N \]
    \[\quad |f(x) - L | < \epsilon.\]

    Hence we write $\lim_{x \rightarrow \infty} f(x) = L$.
\end{definition}

\begin{definition}
    Let $D \subseteq \mathbb{R}$, $c \in \mathbb{R}$ be a cluster point of $D$ and $f: D \rightarrow \mathbb{R}$

    \begin{itemize}
      \item We say $f$ diverges to $\infty$ as $x \rightarrow c$ if $\forall M \in \mathbb{R}$ $\exists \delta >0$ such that $\forall x \in D$ with $0 < | x - c| < \delta$ we have $f(x) > M$. In this case, we write $\lim_{x \rightarrow c} f(x) = \infty$.
      \item We say $f$ diverges to $ - \infty$ as $x \rightarrow c$ if $\forall M \in \mathbb{R}$ $\exists \delta >0$ such that $\forall x \in D$ with $0 < | x - c| < \delta$ we have $f(x) < M$. In this case, we write $\lim_{x \rightarrow c} f(x) = - \infty$.
    \end{itemize}
\end{definition}
